\documentclass[11pt,a4paper]{article}
\usepackage[margin=25mm,headheight=15pt]{geometry}
\usepackage{amsmath,amssymb,amsthm,mathtools,booktabs,array,longtable,microtype}
\usepackage[T1]{fontenc}
\usepackage{lmodern,xcolor,xurl,hyperref,fancyhdr}
\hypersetup{colorlinks=true,linkcolor=blue!45!black,citecolor=blue!45!black,urlcolor=blue!45!black,
 pdftitle={Flow-polynomial obstructions to the random-cluster bunkbed inequality},
 pdfauthor={Anonymous},pdfsubject={Partial resolution; exact certificates; not a complete classification}}
\pagestyle{fancy}\fancyhf{}
\fancyhead[L]{\small Evidence Press\enspace /\enspace Mathematics candidate}
\fancyhead[R]{\small v0.3.1 candidate}
\fancyfoot[L]{\small 30 September 2026}\fancyfoot[R]{\small\thepage}
\setlength{\parindent}{0pt}\setlength{\parskip}{5pt}\setlength{\emergencystretch}{1em}
\renewcommand{\headrulewidth}{0.3pt}
\newtheorem{theorem}{Theorem}[section]
\newtheorem{lemma}[theorem]{Lemma}
\newtheorem{proposition}[theorem]{Proposition}
\newtheorem{corollary}[theorem]{Corollary}
\theoremstyle{remark}\newtheorem{remark}[theorem]{Remark}
\DeclareMathOperator{\tr}{tr}\DeclareMathOperator{\diag}{diag}
\newcommand{\one}{\mathbf 1}\newcommand{\N}{\mathcal N}
\newcommand{\R}{\mathcal R}\newcommand{\CC}{\mathcal C}
\newcommand{\bb}{\square K_2}
\title{\vspace{-10mm}\textbf{Flow-polynomial obstructions to the}\\
\textbf{random-cluster bunkbed inequality}\\[3mm]
{\large An exact homogeneous parameter window\\around the Ising point}}
\author{Anonymous}\date{30 September 2026}
\begin{document}
\maketitle\vspace{-7mm}
\begin{abstract}
We connect random-cluster bunkbed counterexamples to negative evaluations of Eulerian flow polynomials and prove a homogeneous graph-transfer theorem for every real $q>1$. We then determine a whole window of universal validity: for each fixed edge probability $0<p<1$, the only $q$ in $[\alpha,691/250]$ for which the bunkbed inequality holds on every finite graph is $q=2$, where $\alpha=0.787422865474633\ldots$ is the unique real root of $q^5-7q^4+19q^3-28q^2+26q-10$. The upper exclusion combines an analytic complete-bipartite family with exact circulant flow certificates. Below one, a second spectral limit reduces all-$p$ failure to nine polynomial positivity certificates. At $q=1$, a Jordan-mode calculation supplies a separate direct construction. Every negative conclusion concerns finite connected simple full bunkbeds with the same edge probability on every edge; the graph may depend on both parameters. We provide exact symbolic and finite checks, smaller explicit graph recipes, and optimisation-safe verification with deliberate corruption tests. The full parameter classification remains unresolved.
\end{abstract}
\vspace{1mm}
\noindent\fcolorbox{black!35}{black!3}{\parbox{0.965\linewidth}{\small
\textbf{Scope and assurance.} This is a research candidate, not a complete solution of Ayyer--Linusson--Ravichandran Problem 1.4. The manuscript supplies proofs and reproducible exact computations; it is not externally peer reviewed or formally verified. Its novelty has not been established by an exhaustive priority audit.}}

\section{The question and the results}
For a finite simple graph $G=(V,E)$, its full bunkbed is $G\bb$. Write $u_i=(u,i)$, $i\in\{0,1\}$. At $q>0$ and $0<p<1$, set $x=p/(1-p)$. The random-cluster probability of an edge set $A$ is proportional to $q^{k(A)}x^{|A|}$, where isolated vertices count towards $k(A)$. Define
\begin{equation}\label{eq:signed}
\N_G(q,x)=\sum_{A\subseteq E(G\bb)}q^{k(A)}x^{|A|}
\big(\one_{u_0\leftrightarrow v_0}-\one_{u_0\leftrightarrow v_1}\big).
\end{equation}
The partition function is strictly positive. Thus the desired bunkbed inequality holds precisely when $\N_G(q,x)\ge0$.
For each fixed $0<p<1$, define
\[
\mathcal U_p=\{q>0:\N_G(q,p/(1-p))\ge0\text{ for every finite simple }G\text{ and every }u,v\in V(G)\}.
\]


Ayyer, Linusson, and Ravichandran formulate the full parameter question as Problem 1.4, record the positive result at $q=2$, and report counterexamples for $0.56<q<1.43$ \cite{ALR}. Their construction builds on Hollom and on the graph counterexample of Gladkov, Pak, and Zimin \cite{GPZ,Hollom}. We do not assume that the set of valid $q$ is an interval.

GPZ explicitly ask for the least universally valid $q>1$ in Section~8.5 of their published paper and note that validity for all $q\ge2$ was undecided \cite{GPZ}. We settle the least-$q$ question and exclude the latter possibility. Their 8 October 2024 author version, footnote~3, also records Sokal's suggestion of validity at integers $q\ge2$ and failure at nonintegers $q>2$ \cite{GPZauthor}. Our result covers only part of the proposed noninteger range and leaves integer $q\ge3$ untouched.

\paragraph{Comparison of models and quantifiers.}
Hollom's hypergraph counterexample and GPZ's ordinary-graph disproof at $q=1$ are antecedents, not claims of this paper. ALR's Appendix~A starts with fixed posts and unequal edge probabilities $1/100$ and $99/100$, reports the range $0.56<q<1.43$, then discusses full bunkbeds and series--parallel conversion to unweighted examples. Thus prior work is not restricted to weighted or conditioned-post models. Its displayed numerical construction does not state a counterexample at every prescribed common $p$. Here every $p\in(0,1)$ is prescribed first, and for every $q\in[\alpha,2)\cup(2,691/250]$ the construction supplies a finite simple full bunkbed with that same $p$ on all edges; the graph may depend on $(p,q)$. In particular, $\alpha>0.56$ is \emph{not} an improved lower endpoint. The below-one contribution is a direct all-$p$ component of a unified construction; the general Eulerian-flow transfer and the consequences above one are separate structural contributions.

All asymptotic constants below may depend on the fixed word and fixed parameters $q,x$, unless a narrower dependence is stated. No uniformity as $q$ approaches $1$ or $2$, or as $x$ approaches $0$ or infinity, is asserted.


For a multigraph $\Gamma$ on $t$ vertices and $m$ edges, use the flow-polynomial convention
\begin{equation}\label{eq:flow}
F_\Gamma(q)=\sum_{B\subseteq E(\Gamma)}(-1)^{m-|B|}q^{|B|-t+k(B)}.
\end{equation}
Our central transfer result is the following.
\begin{theorem}[Eulerian-flow transfer]\label{thm:transfer}
Let $\Gamma$ be a connected loopless Eulerian multigraph and let $q>1$ satisfy $F_\Gamma(q)<0$. For every $p\in(0,1)$, there is a finite connected simple graph $G$, and vertices $u,v\in V(G)$, for which the homogeneous random-cluster measure on the full bunkbed $G\bb$ satisfies
\[
\mu_{p,q}(u_0\leftrightarrow v_0)<\mu_{p,q}(u_0\leftrightarrow v_1).
\]
The construction is explicit from an Euler circuit, a sufficiently long uniform fan, and sufficiently many pendant vertices at designated posts.
\end{theorem}
\begin{corollary}[Exact local classification]\label{cor:main}
Let $\alpha$ be the unique real root of
\[
g(q)=q^5-7q^4+19q^3-28q^2+26q-10.
\]
For every fixed $p\in(0,1)$,
\begin{equation}\label{eq:main}
\boxed{\qquad \mathcal U_p\cap[\alpha,691/250]=\{2\}.\qquad}
\end{equation}
Equivalently, for each $q\in[\alpha,2)\cup(2,691/250]$ and every $p\in(0,1)$, a finite connected simple full bunkbed violates the inequality. The least universally valid $q>1$ is exactly $2$, and universal validity is not upward-closed in $q$.
\end{corollary}
The graph depends on both parameters. This does not contradict positive results near $p=1$ for each fixed graph \cite{ALR}. At $p=0$ or $p=1$ the inequality holds directly for every $q>0$; these endpoints are excluded from \eqref{eq:main}.

\textbf{Proof dependencies.} The word identity and finite real-$q$ representation yield the Eulerian-flow transfer after exact post amplification. Thickened triangles handle $(1,2)$; complete bipartite graphs and a circulant certificate handle $(2,691/250]$. A different spectral limit and Bernstein certificates handle $[\alpha,1)$, and a separate Jordan calculation handles $q=1$. H\"aggstr\"om's credited Ising theorem supplies the sole positive point. The exact calculations certify finite identities and signs; they do not replace the written limit and amplification arguments.

\section{A word identity}
\subsection{The auxiliary hypergraph}
Let $w=(c_1,\ldots,c_m)$ be a nonempty word using exactly $t$ labels, called posts. Introduce distinct path vertices $z_0,\ldots,z_m$, disjoint from the posts, and hyperedges
\[
h_i=\{z_{i-1},z_i,c_i\}.
\]
Take two path copies and identify the two copies of every post. The two copies of $h_i$ each have absent weight $1$ and present weight $\lambda_i$. A present hyperedge connects its three vertices. Define the signed random-cluster numerator $\N_w(q;\lambda)$ between $z_{0,0}$ and $z_{m,0},z_{m,1}$ by the analogue of \eqref{eq:signed}.

Let $\Gamma_w$ have the posts as vertices and the cyclic transition edges
\[
(c_i,c_{i+1}),\qquad 1\le i\le m,\quad c_{m+1}=c_1.
\]
It is connected and Eulerian; it may have loops or parallel edges.
\begin{theorem}[Exact word identity]\label{thm:word}
As a polynomial identity, with the quotient at $q=1$ interpreted removably,
\begin{equation}\label{eq:word}
\boxed{\displaystyle
\N_w(q;\lambda)=q^{t+2}\left(\prod_{i=1}^m\lambda_i\right)
\frac{F_{\Gamma_w}(q)}{q-1}.}
\end{equation}
\end{theorem}
\begin{proof}
First take integer $q\ge2$ and fix the post colours $s_c\in[q]$. Expanding each factor $1+\lambda_i\one_{\text{all three colours agree}}$ gives the random-cluster weight: every connected component can be coloured in $q$ ways. Put $J=\one\one^\top$ and
\[
M_{s,i}=J+\lambda_i e_se_s^\top,\qquad
K=\prod_{i=1}^m M_{s_{c_i},i},\qquad z=\one^\top K\one.
\]
For two layers conditional on the post colours, the unnormalized difference of same-colour events is
\begin{equation}\label{eq:S}
S=z\tr K-\one^\top K^2\one.
\end{equation}
Indeed, the within-layer term is $z\tr K$, whereas the cross-layer term pairs the first-layer initial marginal with the second-layer final marginal. On averaging all colours, the difference of connectivity indicators is $q/(q-1)$ times the difference of equality indicators.

Write $u_s=\one\wedge e_s$ in the second exterior power. Direct expansion gives
\[
\bigwedge^2 M_{s,i}=\lambda_i u_su_s^\top,
\qquad \langle u_s,u_a\rangle=q\delta_{sa}-1.
\]
For any $K$, the expression in \eqref{eq:S} equals
$\sum_a\langle u_a,(\bigwedge^2K)u_a\rangle$. Since
$\sum_a u_au_a^\top$ is $q$ times the projection onto
$\one\wedge\one^\perp$, and the product here is supported on that subspace,
\[
S=q\left(\prod_i\lambda_i\right)
\prod_{i=1}^m(q\delta_{s_{c_i},s_{c_{i+1}}}-1).
\]
Finally, expanding the product and summing over all post colours gives
\begin{align*}
\sum_{s\in[q]^t}\prod_{ab\in E(\Gamma_w)}(q\delta_{s_a,s_b}-1)
&=\sum_{B\subseteq E(\Gamma_w)}(-1)^{m-|B|}q^{|B|+k(B)}\\
&=q^tF_{\Gamma_w}(q).
\end{align*}
Multiplying by $q/(q-1)$ proves \eqref{eq:word} for every integer $q\ge2$. Both sides, after multiplication by $q-1$, are polynomials, so the identity holds identically. In particular the apparent pole is removable.
\end{proof}
For $w=ABCACB$ or $ABCABC$, the transition graph is the doubled triangle and
\[
F_{2C_3}(q)=(q-1)(q^3-5q^2+10q-7).
\]
This recovers the cubic occurring in the ALR auxiliary calculation. The cubic and its real root near $1.430159709$ already occur in the flow-polynomial literature \cite{Dong}; the cubic itself is not a new result. The contribution here is the general word identity and its homogeneous graph transfer.

\section{Uniform fans and real-$q$ continuation}
\subsection{The ordinary graph}
Take a loopless Eulerian transition graph and an Euler word $w$ of length $m$. Replace the triple $\{a,b,c\}$ at position $i$ by a fan: a rim path
\[
b=y_0-y_1-\cdots-y_L=c
\]
and all $L+1$ spokes $ay_j$. Give every edge the same activity $x>0$. Consecutive fans share the appropriate rim endpoint and no internal rim vertices. Their post labels are distinct at a common endpoint because the transition graph has no loops. Therefore the resulting base graph $G_{w,L}$ is simple, and
\begin{equation}\label{eq:corecounts}
n=t+mL+1,\qquad e=m(2L+1).
\end{equation}
Deleting the $t$ posts leaves a single path whose endpoints are $u,v$.

For the moment, wire the two copies of each designated post and omit all other vertical edges. Let $\N_{w,L}^{T}(q,x)$ denote this conditioned-post numerator.

\subsection{The fan matrix and its two dominant eigenvalues}
Set $f=1+x$, $D_s=I+xe_se_s^\top$, and $R=J+xI$. Conditional on the colour $s$ of the post, the one-layer fan transfer is
\begin{equation}\label{eq:fan}
M_{s,L}=D_s(RD_s)^L.
\end{equation}
On $\operatorname{span}\{e_s,\one-e_s\}$, the matrix $RD_s$ is
\[
\begin{pmatrix}f(1+x)&q-1\\ f&q-1+x\end{pmatrix}.
\]
Write its trace, determinant and eigenvalues as
\begin{equation}\label{eq:eigen}
\tau=f(1+x)+q-1+x,\qquad d=fx(q+x),\qquad
\lambda_\pm=\frac{\tau\pm\sqrt{\tau^2-4d}}2.
\end{equation}
For every real $q>1$, both off-diagonal entries are positive and
\begin{equation}\label{eq:ordering}
0<x<\lambda_-<\lambda_+.
\end{equation}
To check the nontrivial comparison, the characteristic polynomial at $x$ equals
$x(q-1)(f-1)>0$, and $x<\tau/2$. Its two positive roots therefore both exceed $x$.
The complementary colour subspace has eigenvalue $x$.

In the genuine integer-colour representation, if
\[
a_L=\frac{fd^L}{q-1},\qquad \R_s=(\one\wedge e_s)(\one\wedge e_s)^\top,
\]
then
\begin{equation}\label{eq:wedge-limit}
a_L^{-1}\bigwedge^2 M_{s,L}\longrightarrow\R_s.
\end{equation}
The determinant on the distinguished plane is $fd^L$; the other exterior modes have relative decay at most a constant times $(x/\lambda_-)^L$.
A statement about integer $q$ alone would not justify using this limit at a noninteger $q$. The next subsection supplies the required finite algebraic continuation.

\subsection{A finite representation for every real $q$}\label{sec:continuation}
Group the $t$ post colours according to their equality partition $\pi$. If $\pi$ has $r$ blocks, the number of ways to assign distinct colours to those blocks is the falling factorial $(q)_r=q(q-1)\cdots(q-r+1)$.
For each partition use $r+1$ coordinates: its $r$ named colours and a common unnamed-colour coordinate. Put
\[
v=(1,\ldots,1)^\top,\qquad w_r=(1,\ldots,1,q-r)^\top,
\qquad J_r=vw_r^\top.
\]
In this subsection replace $J$ by $J_r$ in \eqref{eq:fan}; each post has the named coordinate corresponding to its block. Let
\[
K_\pi=\prod_{i=1}^m M_{\pi(c_i),L},\qquad z_\pi=w_r^\top K_\pi v.
\]
\begin{lemma}[Finite-length algebraic continuation]\label{lem:finite}
For every fixed finite $L$, the conditioned numerator is given, for all real $q\ne1$, by the finite rational identity
\begin{equation}\label{eq:finite}
\N_{w,L}^{T}(q,x)=\frac q{q-1}\sum_{\pi\in\Pi_t}(q)_{|\pi|}
\left[z_\pi\big(\tr K_\pi+(q-|\pi|-1)x^{Lm}\big)
-w_{|\pi|}^\top K_\pi^2v\right].
\end{equation}
The apparent pole at $q=1$ is removable. Values $q-r<0$ and negative falling factorials are algebraic weights, not probability assignments.
\end{lemma}
\begin{proof}
For each integer $q\ge t+1$, group the colour assignments by equality partition. On the unnamed-colour difference subspace, each $D_s$ is the identity and $R$ acts by $x$; consequently $K$ acts by $x^{Lm}$ on a space of dimension $q-r-1$. Its trace contributes exactly the displayed omitted-colour term. All sums over the retained coordinates use the weights $w_r$, giving $z_\pi$ and $w_r^\top K_\pi^2v$. The colour/connectivity conversion contributes $q/(q-1)$. For fixed finite $L$, multiplication by $q-1$ makes both sides polynomials in $q,x$. Agreement at infinitely many integer $q$, for every $x$, proves the polynomial identity. The left side is a finite random-cluster polynomial, so the apparent pole is removable. No large-$L$ inequality is continued from integer values.
\end{proof}

Here is the limit in this representation. The projection onto the distinguished plane is
\[
P_s=e_se_s^\top+\frac{(v-e_s)(w_r-e_s)^\top}{q-1}.
\]
It commutes with $R,D_s$ and is well defined for every $q>1$. Its complementary projection $W_s=I-P_s$ satisfies
\[
M_{s,L}=H_{s,L}+x^LW_s,\qquad H_{s,L}=M_{s,L}P_s.
\]
The second compound $C_2(H_{s,L})$ is exactly $a_L\R_s$, now with
\begin{equation}\label{eq:rankone}
\R_s=(v\wedge e_s)(w_r\wedge e_s)^\top.
\end{equation}
All other terms in $C_2(M_{s,L})/a_L$ tend to zero by \eqref{eq:ordering}. Also
$\|K_\pi\|=O(\lambda_+^{Lm})$, so the extra term
$x^{Lm}z_\pi/a_L^m$ tends to zero.

For clarity, if $C_2(K)$ is indexed by increasing pairs, define the oriented-minor contraction
\begin{equation}\label{eq:contraction}
\CC_w(C_2(K))=\sum_{i,j,k}w_i\det K_{(i,j),(k,j)}
=(w^\top Kv)\tr K-w^\top K^2v,
\end{equation}
where a repeated index gives zero and reversing a pair changes its sign. Expanding each determinant gives
$\sum_{i,j,k}w_i(K_{ik}K_{jj}-K_{ij}K_{jk})$, whose two sums are exactly the two terms on the right. This also explains why $K^2$, not a transpose-based expression, occurs. Formula \eqref{eq:finite} can therefore be passed to the limit using finitely many finite matrices. The contraction of the product of \eqref{eq:rankone} is the same product $q\prod_i(q\delta_{s_{c_i},s_{c_{i+1}}}-1)$ that appears in the word proof. For each fixed partition, this last contraction identity holds for every integer $q\ge t+1$ and hence polynomially for every real $q$. Summing the equality partitions recovers the colour sum.

\begin{proposition}[Uniform fan limit]\label{prop:fanlimit}
For a fixed word with loopless cyclic transition graph, every $q>1$, and every $x>0$,
\begin{equation}\label{eq:limit}
\boxed{\displaystyle
\lim_{L\to\infty}\frac{\N_{w,L}^{T}(q,x)}{a_L^m}
=q^{t+2}\frac{F_{\Gamma_w}(q)}{q-1}.}
\end{equation}
In particular, a negative flow value gives a negative conditioned-post numerator at some finite $L$.
\end{proposition}
The proof above establishes the limit for real $q$, rather than continuing an inequality from the integers. Appendix~\ref{app:bound} gives an explicit rational bound on its error.

\section{From conditioned posts to the full bunkbed}\label{sec:posts}
\begin{lemma}[Exact post amplification]\label{lem:amplify}
Let a finite simple base graph have a designated post set $T$, distinct vertices $u,v\notin T$, and suppose deleting $T$ leaves a single path from $u$ to $v$ containing all nonpost vertices. Suppose its conditioned-post numerator at $q>0,x>0$ is negative. Attaching sufficiently many pendant base vertices to each post gives a negative numerator on the full bunkbed, with every edge still having activity $x$.
\end{lemma}
\begin{proof}
In the full core, temporarily give the $t=|T|$ designated vertical edges activity $y$, leaving every other edge at activity $x$. Its numerator is a polynomial
\[
\N(y)=\sum_{j=0}^t a_jy^j.
\]
Its leading coefficient is exactly
\begin{equation}\label{eq:leading}
a_t=\N^T.
\end{equation}
To see this, fix all designated verticals open. For a configuration with additional open nonpost verticals, choose the first such vertex along the ordered path from $u$ to $v$. After contracting its open vertical, interchange the two layers strictly after that vertex, keeping all designated posts fixed. The chosen first vertex and the vertical-edge states do not change, so this defines an involution, not a configuration-dependent unpaired symmetry. The graph has no other nonpost edges, and both copies of every post are already wired; therefore the map preserves the edge weights and component count and exchanges the two target events. If an endpoint is wired, the two events already agree, so its signed contribution is zero. The only surviving contribution has all additional verticals closed, and is precisely the conditioned-post numerator.

A pendant base vertex yields, in the full bunkbed, a three-edge path between the copies of its adjacent post. Summing its two internal vertices gives the exact two-terminal factor
\[
D(1+r\delta),\qquad D=q^2+3qx+3x^2,\qquad r=x^3/D>0.
\]
In random-cluster language this is a positive factor $D$ and an effective edge of activity $r$. With $k$ pendants at a post, the effective vertical activity, including the original vertical edge, is
\begin{equation}\label{eq:yk}
y_k=(1+x)(1+r)^k-1\longrightarrow\infty.
\end{equation}
All eliminated factors are positive. Since $a_t<0$, the numerator is negative for sufficiently large $k$.
\end{proof}
The path-suffix cancellation and post-amplification ideas have antecedents in Hollom and GPZ \cite{Hollom,GPZ}; the assertion here is their exact random-cluster implementation. This is a network reduction, not an independence assumption about random-cluster edges. Applying Lemma~\ref{lem:amplify} to Proposition~\ref{prop:fanlimit} proves Theorem~\ref{thm:transfer}.

For a fully specified finite construction, suppose the core has $n$ vertices, $e$ edges, and $a_t\le-\delta<0$. Put
\begin{equation}\label{eq:Bamp}
B=\max(1,q)^{2n}(1+x)^{2e+n-t}.
\end{equation}
Then $|a_j|\le\binom tj B$. For $y\ge1$,
\[
\N(y)\le-\delta y^t+(2^t-1)By^{t-1}.
\]
It is sufficient that $y>\max(1,2^tB/\delta)$. An entirely rational sufficient choice is
\begin{equation}\label{eq:ks}
h=\lceil1/r\rceil,\qquad
2^s>\max(2,2^{t+1}B/\delta),\qquad k=hs.
\end{equation}
Bernoulli's inequality gives $(1+r)^h\ge2$, so this choice satisfies the required bound. The final base graph has $n+tk$ vertices and $e+tk$ edges.

\section{Consequences for the parameter question}
\subsection{Every $1<q<2$ fails}
Let $\ell$ be a positive even integer and let $\ell C_3$ denote a triangle with $\ell$ parallel edges on each side. An Euler word is $(ABC)^\ell$. With $z=q-1$, expansion over the three post colours, or diagonalization of the three-cycle colour transfer, gives
\begin{equation}\label{eq:thick}
F_{\ell C_3}(q)=q^{-3}\left[(z^\ell+z)^3+z(z^\ell-1)^3\right],
\qquad
\frac{F_{\ell C_3}(q)}{q-1}
=\frac{z^{3\ell-1}+3z^\ell+z-1}{q^2}.
\end{equation}
These are identities for the polynomial defined in \eqref{eq:flow}; the displayed denominators cancel. For example, in the colour sum each edge bundle contributes $z^\ell$ when its endpoints have equal colours and $1$ otherwise. The matrix of these weights has eigenvalues $z^\ell+q-1$ and $z^\ell-1$, of multiplicities $1$ and $q-1$; its cubic trace proves the first formula for integer $q$, hence identically. Expanding it gives the second formula.

When $1<q<2$, we have $0<z<1$. Choose an even $\ell\ge2$ such that
\[
4z^\ell<1-z.
\]
Since $z^{3\ell-1}\le z^\ell$, the last numerator in \eqref{eq:thick} is negative. Theorem~\ref{thm:transfer} applies at every prescribed $p\in(0,1)$. For $q=3/2$, one can take $\ell=4$ and obtains
\[
F_{4C_3}(3/2)=-71/1024.
\]
The construction has only three designated posts, although its size grows as $q$ approaches $2$.

\subsection{The known positive point and the exact minimum}
\begin{proposition}[H\"aggstr\"om's Ising point]\label{prop:ising}
At $q=2$ the bunkbed inequality holds on every finite graph for every $p\in[0,1]$.
\end{proposition}
\begin{proof}[Proof using the standard Ising association inequality]
This is the known theorem of H\"aggstr\"om \cite{Haggstrom}; the following argument records the relevant hypotheses. Condition on the open vertical-edge set and contract it. In the Ising representation, further condition on the spins of the resulting posts. The unpinned spins in the two layers are independent copies of the same ferromagnetic Ising model with fixed boundary spins. Arbitrary signs of the induced boundary fields do not affect the FKG lattice condition \cite{FKG}. Thus the conditional covariance of the endpoint spins in one layer is nonnegative. The cross-layer spin product has conditional expectation equal to the product of their means. Averaging the resulting same-layer minus cross-layer inequality and using the random-cluster/Ising colour expansion proves the claim. If an endpoint is pinned, its conditional covariance is zero. Finally average over the vertical-edge set. The endpoints $p=0,1$ follow directly.
\end{proof}
Every $q\in(1,2)$ fails, whereas $q=2$ succeeds. This proves the minimum assertion in Corollary~\ref{cor:main}; it is a minimum, not merely an infimum.

\subsection{An obstruction above $2$}\label{sec:above}
Let $\Gamma_*$ be the simple graph on $\{0,\ldots,11\}$ specified by the following closed Euler word:
\begin{equation}\label{eq:wordstar}
\begin{split}
w_*={}&(0,3,8,1,5,6,2,10,5,4,9,3,\\
       &6,11,9,1,10,7,4,0,7,8,2,11).
\end{split}
\end{equation}
Each cyclically consecutive pair is one unordered edge. The complete edge list is also supplied in \texttt{flow\_witness.json}.
The closing transition is $11\to0$. The graph is connected and $4$-regular. Direct exact enumeration gives $F_{\Gamma_*}(q)=(q-1)P(q)$, where
\begin{equation}\label{eq:poly}
\begin{split}
P(q)={}&q^{12}-23q^{11}+253q^{10}-1759q^9+8615q^8-31366q^7\\
 &+87285q^6-187746q^5+311174q^4-389370q^3\\
 &+350419q^2-203812q+57709.
\end{split}
\end{equation}
In particular,
\begin{equation}\label{eq:negativevalue}
F_{\Gamma_*}(21/10)=-\frac{20868388946779}{10^{13}}<0,
\qquad
F_{\Gamma_*}(9/4)=-\frac{57285175}{67108864}<0.
\end{equation}
Endpoint values alone do not prove negativity between them. Appendix~\ref{app:bernstein} gives all thirteen strictly negative Bernstein coefficients of $P$ on $[21/10,9/4]$, proving negativity on the entire closed interval. Theorem~\ref{thm:transfer} now proves its failure assertion at every $p\in(0,1)$. Since $2$ itself is a universal positive point, universal validity is not upward-closed.

\subsection{An analytic interval immediately above the Ising point}\label{sec:hub}
Let $b$ be a positive even integer. The complete bipartite graph $K_{4,b}$ is connected, simple and Eulerian. Writing $z=q-1$, its flow polynomial satisfies
\begin{equation}\label{eq:hub}
\begin{split}
q^{b+3}F_{K_{4,b}}(q)={}&(z^4+z)^b+4z(-1)^b(z^3+1)^b\\
 &+3z(2z^2+z-1)^b\\
 &+6z(z-1)(z^2-z-2)^b\\
 &+z(z-1)(z-2)(-3q)^b.
\end{split}
\end{equation}
To derive this identity, first take integer $q\ge4$ and sum the colour-product expression in the word proof over the four vertices on the small side. Their equality-partition types are $4$, $3+1$, $2+2$, $2+1+1$, and $1+1+1+1$, with multiplicities $1,4,3,6,1$. Summing the colour of one vertex on the other side gives, respectively, the five bases displayed in \eqref{eq:hub}. The falling-factorial weights give its coefficients. Polynomial continuation proves the formula at real $q$; this is an identity, not continuation of a sign assertion.

Let $q_c$ be the unique real root of
\[
q^3-4q^2+6q-6=0,
\qquad q_c=2.574743073887\ldots.
\]
Uniqueness follows because the derivative is $3(q-4/3)^2+2/3>0$. For $2<q<q_c$, the last coefficient in \eqref{eq:hub} is negative, while the other coefficients are positive. Its absolute base $3q$ strictly exceeds all other absolute bases. Indeed,
\[
\begin{array}{ll}
3q-(z^4+z)=-q(q^3-4q^2+6q-6)>0,&
3q-(z^3+1)=q^2(3-q)>0,\\
3q-(2z^2+z-1)=2q(3-q)>0,&
|z^2-z-2|=q(3-q)<3q.
\end{array}
\]
Put $\rho=\max\{z^4+z,z^3+1,2z^2+z-1,|z^2-z-2|\}/(3q)<1$. Any sufficiently large even $b$ with
\[
(1+z+6z^2)\rho^b<z(z-1)(2-z)
\]
makes the flow polynomial negative. Theorem~\ref{thm:transfer} therefore excludes every $2<q<q_c$ at every prescribed $p\in(0,1)$. This argument makes no claim that $K_{4,b}$ is optimal among complete bipartite or other Eulerian families.

\subsection{The certified extension to $691/250$}\label{sec:circulant}
For distinct steps $a,b$, let $C_n(a,b)$ be the simple circulant on $\mathbb Z/n\mathbb Z$ with unordered edges $\{i,i+a\}$ and $\{i,i+b\}$. In particular $C_{32}(1,5)$ is connected and $4$-regular, with $32$ vertices and $64$ edges. Its exactly reconstructed polynomial satisfies
\begin{equation}\label{eq:c32}
F_{C_{32}(1,5)}(q)<0\quad\hbox{for}\quad
\frac{20000001}{10000000}\le q\le\frac{691}{250}.
\end{equation}
The certificate supplies the full integer coefficient vector, divides it exactly by $q-1$, checks negative values at both rational endpoints, and uses an exact Sturm sequence to count zero roots in the interval. The frontier computation reconstructs the vector from the edge list rather than trusting a stored polynomial. The accompanying structural check verifies connectedness, simplicity and every vertex degree. No approximate root is a sign certificate.

Since $20000001/10000000<q_c$, \eqref{eq:c32} overlaps the analytic interval of Section~\ref{sec:hub}. Hence their union excludes the whole interval
\[
\boxed{\quad 2<q\le691/250=2.764\quad}
\]
at every prescribed interior $p$. The following selected witnesses illustrate the certified progression; the package checks sixteen graphs in total.
\begin{center}
\begin{tabular}{lrrl}
\toprule
Auxiliary graph & Vertices & Edges & Certified negative interval\\
\midrule
$C_{10}(1,4)$ &10&20&$[2.0416727,\ 2.2898]$\\
$C_{16}(2,3)$ &16&32&$[2.0006061,\ 2.5655]$\\
$C_{24}(3,4)$ &24&48&$[2.0000026,\ 2.7102]$\\
$C_{28}(1,5)$ &28&56&$[2.0000003,\ 2.745]$\\
$C_{32}(1,5)$ &32&64&$[2.0000001,\ 2.764]$\\
\bottomrule
\end{tabular}
\end{center}
Every terminating decimal in this table denotes its exact rational value. These vertex counts belong to auxiliary flow graphs, not to the much larger ordinary bunkbed counterexamples. No maximality of the endpoint $2.764$ is asserted.
\section{A universal obstruction below $q=1$}\label{sec:below}
The order of the fan eigenvalues changes at $q=1$. For $0<q<1$ the characteristic polynomial at $x$ is negative, so
\[
0<\lambda_-<x<\lambda_+.
\]
Consequently Proposition~\ref{prop:fanlimit} does not apply. A different normalisation retains the mixed mode between the top eigenvalue and the complementary eigenvalue $x$.

\subsection{The second fan limit}
In a reduced representation write $\Pi_s=I-P_s$, and let $A_s^{+}$ be the rank-one coefficient of $\lambda_+^L$ in $M_{s,L}$. Define
\[
\Lambda_s=C_2(A_s^{+}+\Pi_s)-C_2(A_s^{+})-C_2(\Pi_s).
\]
For a post partition $\pi$, put $Q_\pi=\prod_i A_{\pi(c_i)}^{+}$. The finite identity \eqref{eq:finite} and the strict spectral ordering give the following limit.
\begin{proposition}[The subunit-$q$ fan limit]\label{prop:belowlimit}
For each fixed $0<q<1$, $x>0$, and word $w$ with $m$ letters,
\begin{equation}\label{eq:belowlimit}
\begin{split}
G_w(q,x):={}&\lim_{L\to\infty}\frac{\N^T_{w,L}(q,x)}{(\lambda_+x)^{Lm}}\\
={}&\frac q{q-1}\sum_{\pi\in\Pi_t}(q)_{|\pi|}
\left[\CC_{w_r}\!\left(\prod_i\Lambda_{\pi(c_i)}\right)
+(q-r-1)w_r^\top Q_\pi v\right],\quad r=|\pi|.
\end{split}
\end{equation}
If $G_w(q,x)<0$ and the cyclic transition graph is loopless, some finite homogeneous fan core has negative numerator; Lemma~\ref{lem:amplify} then gives a full homogeneous bunkbed counterexample.
\end{proposition}
\begin{proof}
At fixed $q,x$, the spectral decomposition is
$M_{s,L}=\lambda_+^L A_s^{+}+\lambda_-^L A_s^{-}+x^L\Pi_s$.
The second compound of the top rank-one term vanishes. After division by $(\lambda_+x)^L$, its mixed compound with $\Pi_s$ survives. Every other term decays, since $\lambda_-/x<1$ and $x/\lambda_+<1$. The omitted-colour term in \eqref{eq:finite} now survives: $K_\pi/\lambda_+^{Lm}\to Q_\pi$. There are finitely many partitions, so these are legitimate termwise limits of a finite real-$q$ identity. All normalising factors are positive.
\end{proof}

\subsection{An explicit polynomial for the doubled triangle}
Take $w=ABCABC$. Put $\lambda=\lambda_+$, $f=1+x$, and
\begin{equation}\label{eq:hs}
h=x+\frac{fx^2}{\lambda-fx},\qquad
s=q-1+\frac{(1+h)^2}{f}.
\end{equation}
The characteristic polynomial at $fx$ equals $-fx^2<0$, so $\lambda>fx$, $h>x>0$, and $s>q+x>0$.
The right top eigenvector of $RD_s$ can be chosen as
$\ell_s=v+x^2e_s/(\lambda-fx)$. The weighted left eigenvector is
$\ell_s^\top\operatorname{diag}(w_r)D_s/s$.
It follows that
\[
A_s^{+}=\frac{(v+he_s)(w_r+he_s)^\top}{s}.
\]
This formula uses an algebraic bilinear form when $q-r<0$; it does not introduce a signed probability measure.

\begin{lemma}[Doubled-triangle polynomial]\label{lem:belowpoly}
There is an explicit polynomial $P(q,h)=\sum_{j=0}^{10}c_j(q)h^j$, given in Appendix~\ref{app:belowcoeff}, for which
\begin{equation}\label{eq:belowpoly}
\boxed{\quad
G_{ABCABC}(q,x)=
\frac{q^2(h+q)^2(q-2)}{s^6(q-1)^6}\,P(q,h).
\quad}
\end{equation}
Its first two coefficients are
\[
c_0(q)=q^7g(q),\qquad c_1(q)=10q^6g(q),
\quad g(q)=q^5-7q^4+19q^3-28q^2+26q-10.
\]
\end{lemma}
\begin{proof}[Finite algebraic proof]
For each of the five equality partitions $000,001,010,011,012$, use its own reduced matrices and put
\[
A_s=(v+he_s)(w_r+he_s)^\top,\qquad
B_s=(q-1)(I-e_se_s^\top)-(v-e_s)(w_r-e_s)^\top.
\]
Then $A_s^{+}=A_s/s$ and $\Pi_s=B_s/(q-1)$. Let
$L_s=C_2(A_s+B_s)-C_2(A_s)-C_2(B_s)$.
Thus $\Lambda_s=L_s/[s(q-1)]$. Direct polynomial multiplication gives
\begin{equation}\label{eq:belowidentity}
\begin{split}
&\sum_{\pi\in\Pi_3}(q)_r\left[
\CC_{w_r}\!\left(\prod_{i=1}^{6}L_{\pi(c_i)}\right)
+(q-1)^6(q-r-1)w_r^\top\left(\prod_{i=1}^{6}A_{\pi(c_i)}\right)v\right]\\
&\hspace{25mm}=q(h+q)^2(q-2)(q-1)P(q,h).
\end{split}
\end{equation}
Substitution into \eqref{eq:belowlimit} proves the claim. This small symbolic identity is reconstructed from the five matrices by \texttt{verify\_below\_universal.py}, and its coefficient vector is printed in Appendix~\ref{app:belowcoeff}. The three two-block partitions must be evaluated separately: their ordered contractions cannot be replaced by three times one representative.
An additional standard-library checker, \texttt{code/check\_below\_identity.py}, evaluates both sides at the $18\times13$ integer grid $q=3,\ldots,20$, $h=0,\ldots,12$. To justify this certificate, each entry of $A_s$ has bidegree at most $(1,2)$ and each entry of $B_s$ at most $(1,0)$. The mixed minor $L_s$ has bidegree at most $(2,2)$. Six factors, the contraction weight, and the falling factorial give degree at most $12+1+3=16$ in $q$ for the first term. The second term has degree at most $6+1+1+6+3=17$. Both have degree at most $12$ in $h$. The printed right side has the same upper bounds (in fact degree at most $17$ in $q$ and $12$ in $h$). Successive univariate interpolation therefore makes equality on this grid an exact identity certificate, not a numerical sign test.
\end{proof}

\subsection{A whole interval, at every edge probability}
Let $\alpha$ denote the unique real root of $g$. Exact root isolation gives
\begin{equation}\label{eq:alpha}
\frac{787422865474633}{10^{15}}<\alpha<
\frac{787422865474634}{10^{15}},
\qquad \alpha=0.7874228654746333597\ldots.
\end{equation}
\begin{theorem}[All-$p$ failure below one]\label{thm:below}
For every $\alpha\le q<1$ and every $0<p<1$, there is a finite connected simple homogeneous full-bunkbed counterexample.
\end{theorem}
\begin{proof}
For $2\le j\le10$, all Bernstein coefficients of $c_j(q)$ on $[787/1000,1]$ are strictly positive. The certificate records the rational coefficients and checks the reverse change of basis. Since $\alpha>787/1000$, this proves $c_j(q)>0$ for $\alpha\le q<1$.
The quintic has exactly one real root and positive leading coefficient; its certified bracket has $g$ negative at the left endpoint and positive at the right. Hence $g(q)\ge0$ for $q\ge\alpha$, so $c_0,c_1\ge0$ there. Since $h>0$, $P(q,h)>0$, including at $q=\alpha$, where its first two coefficients vanish. The prefactor in \eqref{eq:belowpoly} is strictly negative for $0<q<1$. Proposition~\ref{prop:belowlimit} and Lemma~\ref{lem:amplify} finish the proof.
\end{proof}
In particular, the simpler rational interval $4/5\le q<1$ needs no algebraic endpoint: all eleven coefficient polynomials have strictly positive Bernstein coefficients on $[4/5,1]$.

The endpoint $\alpha$ is sharp only for the \emph{all-$p$ negativity of this fixed limiting test}. For a fixed $0<q<\alpha$, as $x\downarrow0$ we have $h\to0$, $s\to q$, and $c_0(q)<0$. Therefore $G_{ABCABC}(q,x)>0$ for sufficiently small $x$. This says nothing about other words or other graph constructions, and does not prove the bunkbed inequality below $\alpha$.
\section{The boundary point $q=1$}\label{sec:percolation}
Neither preceding spectral limit can simply be evaluated at $q=1$. A polynomial factor in the fan length replaces the coinciding exponential modes. The following gives a direct homogeneous all-$p$ construction within the same framework. It does not claim a new disproof of Bernoulli bunkbed percolation, which is due to GPZ \cite{GPZ}.

\begin{theorem}[The percolation fan limit]\label{thm:qone}
For $w=ABCABC$ and every fixed $x>0$, set $f=1+x$ and $\kappa=x/(x^2+x+1)$. Then
\begin{equation}\label{eq:qonelimit}
\boxed{\quad
\lim_{L\to\infty}
\frac{\N^T_{w,L}(1,x)}{(f\kappa L)^6(f^2x)^{6L}}=-1.
\quad}
\end{equation}
Consequently, at $q=1$, every prescribed $0<p<1$ admits a finite connected simple homogeneous full-bunkbed counterexample.
\end{theorem}
\begin{proof}
Fix a partition with $r$ blocks and put $q=1$ in its reduced matrices. Write
$T_s=(vw_r^\top+xI)D_s$ and $U_s=T_s-xI$. Its minimal polynomial divides
$(X-f^2)(X-x)^2$. Since $f^2-x=x^2+x+1>0$, the matrices
\[
E_s=\frac{U_s^2}{(f^2-x)^2},\quad
A_s=D_sE_s,\quad B_s=D_s(I-E_s),\quad
C_s=x^{-1}D_sU_s(I-E_s)
\]
satisfy $E_s^2=E_s$, $T_sE_s=f^2E_s$, and $U_s^2(I-E_s)=0$. Therefore
\[
M_{s,L}=f^{2L}A_s+x^L(B_s+LC_s).
\]
The top matrix $A_s$ has rank one. Its mixed compound with $C_s$ is exactly
\[
C_2(A_s+C_s)-C_2(A_s)-C_2(C_s)
=f\kappa\,(v\wedge e_s)(w_r\wedge e_s)^\top
=f\kappa\R_s.
\]
These are identities of rational functions of $x$; the accompanying checker verifies them symbolically, including nilpotence. Hence
\[
\frac{C_2(M_{s,L})}{L(f^2x)^L}\longrightarrow f\kappa\R_s.
\]

To remove the apparent pole in \eqref{eq:finite}, take the derivative of its finite numerator at $q=1$, before passing to the large-$L$ limit. For $r=2,3$, $(q)_r$ has a simple zero at $1$; its derivative is $1,-1$, respectively. Thus only the value of each corresponding bracket at $q=1$ enters. Its omitted-colour term is $O((f^2x)^{6L})$, and disappears after the additional $L^6$ normalisation. The four remaining compound contractions are
\[
\begin{array}{c|rrrr}
\pi&001&010&011&012\\\hline
\CC_{w_r}(\prod_{i=1}^6\R_{\pi(c_i)})&0&0&0&1.
\end{array}
\]
Their weighted sum is $-1$.

For completeness, the one-block partition contributes no term of order $L^6$. In its two-dimensional representation, let $M=M_{0,L}$, $K=M^6$, and $z=w_1^\top K v$. Cayley--Hamilton gives its bracket exactly as
\[
B_1(q)=q\det(M)^6+(q-2)x^{6L}z,
\qquad \det M=f\{fx(q+x)\}^{L}.
\]
At $q=1$, $M$ is triangular, $z=f^{12L+6}$, and $B_1(1)=0$. The top eigenvalue of $RD_0$ remains simple in a neighbourhood of $1$, separated from $x$. Differentiating its spectral representation therefore gives $z'=O(Lf^{12L})$ at $1$; the lower modes and their derivatives are exponentially smaller, with at most polynomial factors in $L$. The displayed determinant formula likewise yields
for some finite $C(x)>0$, independent of the integer $L\ge1$,
\[
|B'_1(1)|\le C(x)(1+L)(f^2x)^{6L},\qquad
\frac{|B'_1(1)|}{L^6(f^2x)^{6L}}
\le C(x)(L^{-6}+L^{-5})\longrightarrow0.
\]
The spectral gap is positive for each fixed $x>0$; its associated constants need not be uniform in $x$. Since the derivative of $qB_1$ at $1$ equals $B'_1(1)$, this proves that the one-block contribution is lower order.

Combining the partitions proves \eqref{eq:qonelimit}. The denominator is positive, so some finite fan has negative numerator; post amplification completes the full graph.
\end{proof}
The coefficientwise differentiation is independently implemented with first-order polynomial jets. Nine finite cases agree exactly with a connectivity-frontier implementation. At $x=1/2,1,2$, the length-$50$ and length-$100$ cores are negative by exact arithmetic. These tests corroborate the formula but do not replace the proof for every real $x>0$.

\section{Finite witnesses and reproducibility}\label{sec:certificates}
\subsection{Explicit graph recipes at $p=1/2$}
At $q=3/2$, the improved core uses $w=(ABC)^4$ and $L=19$. Its full-core post polynomial $\N(y)=a_0+a_1y+a_2y^2+a_3y^3$ is reconstructed exactly by connectivity-frontier dynamic programming; $a_0,a_1,a_2>0>a_3$. With $r=1/(q^2+3q+3)$ and $y_k=2(1+r)^k-1$, exact rational evaluations give
\[
\N(y_{3565})<0\le\N(y_{3564}).
\]
Descartes' rule of signs gives exactly one positive zero, and $y_k$ increases strictly. Thus $3565$ is the minimal equal number of pendants per post for this \emph{fixed core and this construction}, not a globally smallest counterexample. The same calculation at $q=9/10$, $w=(ABC)^2$, $L=21$ gives the minimal equal count $1099$.

The retained conservative $q=21/10$ construction uses $w_*$, $L=1000$, and $1515626$ pendants per post. Here $\tau=61/10$, $d=31/5$, the discriminant is $1241/100>3^2$, and the characteristic polynomial at $5/4$ equals $11/80>0$. Thus $x/\lambda_-<4/5$. Appendix~\ref{app:bound} certifies
\begin{equation}\label{eq:finitegap}
\left|\frac{\N^T_{w_*,1000}(21/10,1)}{a_{1000}^{24}}-C_*\right|<10^{-36},
\quad C_*=-\frac{6154125714915266481205245044409}{10^{26}}<-60000,
\end{equation}
with $a_{1000}=(20/11)(31/5)^{1000}$. Therefore $\delta=(-C_*/2)a_{1000}^{24}$ is a valid bound for amplification; \eqref{eq:ks} gives $h=14$, $s=108259$, $k=1515626$.

\begin{center}
\begin{tabular}{lrrr}
\toprule
&$q=9/10$&$q=3/2$&$q=21/10$\\\midrule
Core base vertices&130&232&24013\\
Core base edges&258&468&48024\\
Pendants per post&1099&3565&1515626\\
Final base vertices&3427&10927&18211525\\
Final base edges&3555&11163&18235536\\
Full bunkbed vertices&6854&21854&36423050\\
Full bunkbed edges&10537&33253&54682597\\\bottomrule
\end{tabular}
\end{center}
The two smaller base edge lists are supplied. The largest graph is represented by a deterministic streaming generator; its full partition function is not directly evaluated. All entries use $|V(G\square K_2)|=2|V(G)|$ and $|E(G\square K_2)|=2|E(G)|+|V(G)|$.

\subsection{Independent checks and their limits}
The twelve-vertex flow polynomial is reconstructed in two different ways. The first enumerates all $2^{24}=16777216$ edge subsets in \eqref{eq:flow}, using a rollback union--find calculation of component counts. The second enumerates all $4213597$ partitions of its twelve vertices. If $h(\pi)$ counts edges whose endpoints lie in the same block, the independent reconstruction is
\begin{equation}\label{eq:partitionflow}
q^{12}F_{\Gamma_*}(q)=\sum_{\pi\in\Pi_{12}}(q)_{|\pi|}
(-1)^{24-h(\pi)}(q-1)^{h(\pi)}.
\end{equation}
Both give exactly the same integer coefficient vector. The Bernstein change of basis is also reconstructed in the reverse direction.

The $q=3/2$ core check uses integer arithmetic with explicit denominator clearing. Its frontier holds the designated posts, the initial vertex, and the two current rim vertices. Closing a component multiplies its weight by $q$. A second implementation uses rational matrices in \eqref{eq:finite}. Four smaller cases, including noninteger $q<1$, are checked against direct edge-subset enumeration as well. Ten hypergraph words, including loops and repeated labels, are checked against \eqref{eq:word}; all surviving multivariate activity exponents are exactly one.

These are algorithmically independent checks, not independent human reviews. The general transfer theorem is justified by the written proof and the explicit bounds, not by numerical testing. The $q=21/10$ certificate checks a sufficient analytic bound; it does not directly evaluate the large graph's partition function. The current verifier uses Python and SymPy for exact polynomial arithmetic, root isolation and symbolic identities; a C++17 compiler is used for the two baseline enumerations. The full mode also reconstructs every developed witness polynomial, including $C_{32}(1,5)$, and both stored full-core cubics. The separate $q=1$ polynomial-jet implementation agrees with a frontier recurrence on nine finite cases. The new below-one formula agrees with a distinct quadratic-field implementation at seven rational parameter pairs.

\section{Scope, priority, and unresolved parts}
Equation~\eqref{eq:main} is an exact local classification, not a complete answer to ALR Problem~1.4. We do not classify all pairs with $0<q<\alpha$ or $q>691/250$. The package retains additional pointwise subunit-$q$ negative limits and finite examples; these are not an all-$p$ theorem below $\alpha$. Nor does a numerical search failure establish universal validity.

No result here settles integer $q\ge3$. Indeed, an Eulerian orientation carrying any constant nonzero element of a group of integer order $q\ge2$ is a nowhere-zero flow. Hence $F_\Gamma(q)\ge q-1>0$ for each Eulerian $\Gamma$, and a negative-flow witness cannot address those integers. This limitation concerns the transfer's sufficient condition, not the truth of the bunkbed inequality there.

The endpoint $691/250$ is a certified rational stopping point, not a claimed optimum; $\alpha$ is a sharp threshold only for the all-$p$ sign of the doubled-triangle limiting test used below one. The graph always may depend on both parameters. No graph-size minimality is claimed beyond the two explicitly fixed-core pendant optimisations.

The original transfer and v0.2 extensions are retained with their provenance. This revision adds the all-$p$ subunit interval, the explicit endpoint polynomial, the separate $q=1$ calculation, and verification repairs. Primary-source retrieval as of 30 September 2026 found no equivalent general Eulerian-flow-to-homogeneous-bunkbed theorem in the retrieved material. The search was targeted, not exhaustive, and cannot certify absence of prior art. The supplied referee report is a model-assisted review of v0.1, not an external human review of the new theorems. External mathematical review remains necessary before treating this candidate as an independently validated publication.

\appendix
\section{A rational error bound for the fan limit}\label{app:bound}
We give the explicit estimate used by \texttt{verify\_asymptotic.py}. Fix $q>1$, $x>0$, $f=1+x$, $t,m\ge1$. Choose positive rational $g,\rho$ with
\[
g\le\sqrt{\tau^2-4d},\qquad x/\lambda_-\le\rho<1.
\]
For rational $q,x$, these inequalities can be certified by polynomial comparisons without approximating a square root. Define
\begin{align*}
S&=\max(q,2t-q),&P_0&=\max\left(1,\frac{2t-q-1}{q-1}\right),&W_0&=1+P_0,\\
C&=\frac{P_0\{2f(S+x)+\tau\}}{g},&
C_0&=(q-1)\left(2CW_0+\frac{W_0^2}{f}\right),\\
\varepsilon_L&=C_0\rho^L,&A_0&=(q-1)P_0+1,&
H_0&=(q-1)(C+W_0/f),\\
U&=q+t+1,&B_t(q)&=\sum_{r=1}^t\left\{\!\begin{matrix}t\\r\end{matrix}\!\right\}|(q)_r|.
\end{align*}
The braces denote Stirling numbers of the second kind.
\begin{lemma}\label{lem:error}
If $\varepsilon_L\le1$, then the absolute error in \eqref{eq:limit} is at most
\begin{equation}\label{eq:error}
E_L=\frac q{q-1}B_t(q)
\left[tS\,mA_0^{m-1}\varepsilon_L
+US H_0^m\rho^{Lm}\right].
\end{equation}
\end{lemma}
\begin{proof}
Use the maximum absolute row-sum norm in each of the reduced representations. For a partition with $r\le t$ blocks, $\|w_r\|_1\le S$, $\|P_s\|_\infty\le P_0$, and $\|W_s\|_\infty\le W_0$. The two spectral projections of $RD_s$ on the distinguished plane are
\[
\frac{(RD_s-\lambda_-I)P_s}{\lambda_+-\lambda_-},\qquad
\frac{(\lambda_+I-RD_s)P_s}{\lambda_+-\lambda_-}.
\]
Since $\|RD_s\|_\infty\le f(S+x)$, their norm sum is at most $C$. Hence
$\|H_{s,L}\|_\infty\le fC\lambda_+^L$.

Expansion of each two-by-two minor gives, for square matrices $A,B$,
\[
\|C_2(B)\|_\infty\le\|B\|_\infty^2,
\quad
\|C_2(A+B)-C_2(A)-C_2(B)\|_\infty
\le2\|A\|_\infty\|B\|_\infty.
\]
Applying these to $M_{s,L}=H_{s,L}+x^LW_s$, and using
$C_2(H_{s,L})=a_L\R_s$, gives
\[
\left\|a_L^{-1}C_2(M_{s,L})-\R_s\right\|_\infty
\le(q-1)\left[2CW_0(x/\lambda_-)^L+
\frac{W_0^2}{f}(x^2/d)^L\right]\le\varepsilon_L.
\]
The rank-one formula \eqref{eq:rankone} gives
$\|\R_s\|_\infty\le(q-1)P_0$. A telescoping product of $m$ factors is therefore within
$mA_0^{m-1}\varepsilon_L$ of the limiting product. The contraction \eqref{eq:contraction} has norm at most $r\|w_r\|_1\le tS$: for each ordered row pair, its oriented column pairs form a subset of the column pairs in one row of the compound matrix.

It remains to bound the omitted-colour contribution. We have
\[
\|M_{s,L}\|_\infty\le f(C+W_0/f)\lambda_+^L,
\quad |q-r-1|\le U,
\quad |z_\pi|\le S\|K_\pi\|_\infty.
\]
Dividing $|(q-r-1)x^{Lm}z_\pi|$ by $a_L^m$ consequently gives at most
$US H_0^m\rho^{Lm}$. Sum the two bounds over partitions with their absolute falling-factorial weights and multiply by $q/(q-1)$ as in \eqref{eq:finite}. This proves \eqref{eq:error}.
\end{proof}
At $q=21/10,x=1,t=12,m=24,L=1000$, the certificate uses $g=3$, $\rho=4/5$. Some intermediate exact constants are
\[
S=219/10,\quad P_0=19,\quad W_0=20,\quad
C=18563/30,\quad C_0=411686/15,\quad H_0=207493/300.
\]
Direct rational comparisons verify $\varepsilon_L<1$, $E_L<10^{-36}$ and $E_L<-C_*/2$. Powers in these comparisons are computed as integers or rational numbers; logarithms and floating-point eigenvalues are not proof dependencies.

\section{The closed-interval sign certificate}\label{app:bernstein}
Put $a=21/10$, $b=9/4$, and $z=(q-a)/(b-a)$. The polynomial in \eqref{eq:poly} has the expansion
\[
P(a+(b-a)z)=\sum_{j=0}^{12}\beta_j\binom{12}{j}z^j(1-z)^{12-j}.
\]
The thirteen Bernstein basis functions are nonnegative on $[0,1]$ and sum to one. Every coefficient in the following table is strictly negative, including the endpoints. It follows that $P(q)<0$ on the entire closed interval, not just on a sampled grid.
\begingroup\renewcommand{\arraystretch}{1.8}
\begin{center}
\begin{tabular}{r@{\qquad}l}
\toprule
$j$ & $\beta_j$\\
\midrule
0 & $-1897126267889/1000000000000$\\
1 & $-1638507048299/800000000000$\\
2 & $-472507171683/220000000000$\\
3 & $-15476196947443/7040000000000$\\
4 & $-3100190903929/1408000000000$\\
5 & $-486611737077/225280000000$\\
6 & $-65422642003/31539200000$\\
7 & $-70138853011/36044800000$\\
8 & $-12799896687/7208960000$\\
9 & $-1803895261/1153433600$\\
10 & $-34375609/26214400$\\
11 & $-85374093/83886080$\\
12 & $-11457035/16777216$\\
\bottomrule
\end{tabular}
\end{center}\endgroup
The checker derives this table from the independently reconstructed flow polynomial, then reconstructs the power coefficients from the table as a second check.

\section{The below-one polynomial and positivity certificates}\label{app:belowcoeff}
The coefficient vector of $P(q,h)=\sum_{j=0}^{10}c_j(q)h^j$ in Lemma~\ref{lem:belowpoly} is as follows. Put $g=q^5-7q^4+19q^3-28q^2+26q-10$.
\begin{align*}
c_0={}&q^7g,\\
c_1={}&10q^6g,\\
c_2={}&5q^5(9q^5-62q^4+165q^3-239q^2+222q-86),\\
c_3={}&40q^4(3q^5-20q^4+51q^3-71q^2+66q-26),\\
c_4={}&2q^3(q^6+100q^5-657q^4+1577q^3-2049q^2+1909q-776),\\
c_5={}&4q^2(3q^6+48q^5-347q^4+783q^3-911q^2+855q-368),\\
c_6={}&-q(q^7-43q^6+734q^4-1741q^3+1723q^2-1748q+864),\\
c_7={}&-4(q^7-23q^6+80q^5-86q^4+39q^3-57q^2-48q+64),\\
c_8={}&2q^6+38q^5-209q^4+431q^3-579q^2+634q-272,\\
c_9={}&2(6q^5-38q^4+111q^3-185q^2+169q-58),\\
c_{10}={}&(q^2-3q+3)(2q^3-9q^2+14q-6).
\end{align*}
For $a=787/1000$, write $c_j(a+(1-a)t)=\sum_{i=0}^{d_j}\beta_{ji}\binom{d_j}{i}t^i(1-t)^{d_j-i}$. The following are strict rational lower bounds for every coefficient in each row. Full rational vectors, rather than only these bounds, are stored in \texttt{below\_universal\_v03.json}.
\begingroup\renewcommand{\arraystretch}{1.8}
\begin{center}\begin{tabular}{rrr}
\toprule
$j$&$d_j$&Certified lower bound for $\min_i\beta_{ji}$\\\midrule
2&10&$\frac{69}{1000}$\\
3&9&$\frac{917}{1000}$\\
4&9&$\frac{171}{40}$\\
5&8&$\frac{10893}{1000}$\\
6&8&$\frac{4249}{250}$\\
7&7&$\frac{2093}{125}$\\
8&6&$\frac{5101}{500}$\\
9&5&$\frac{3519}{1000}$\\
10&5&$\frac{263}{500}$\\
\bottomrule\end{tabular}\end{center}\endgroup
Each lower bound is positive. The verifier reconstructs all coefficients from the five-partition identity, converts them to Bernstein form, and reconstructs the original power coefficients in reverse. Separately, a Sturm count proves that $g$ has exactly one real root and verifies both signs in the rational bracket \eqref{eq:alpha}. The endpoint is included because the strictly positive $h^2,\ldots,h^{10}$ terms remain when $g(\alpha)=0$.

\section{A version-specific source issue and an interface bound}\label{app:source}
ALR Lemma~A.1 as printed in arXiv:2509.18788v1 bounds a projected edge-set measure using the number of common \emph{edges} of two graphs \cite[p.~30]{ALR}. Literally, this is false: on vertices $\{a,b,c\}$ let $G$ have only $ab$, and $H$ have $ac,bc$. There are no common edges. At $q=2$, $p=1/2$, the probability that $ab$ is open is $1/3$ in $G$ and $5/14$ in $G\cup H$, whereas the printed zero-common-edge bound would require equality.

This observation does not disprove ALR's Theorem~1.6. Their three attachment gadgets are compatible with an interface-based bound of the required form. The following precise replacement explains why.
\begin{lemma}[Attachment through a vertex interface]\label{lem:interface}
Let $G,H$ have disjoint edge sets and vertex intersection $S$, where $b=|S|\ge1$. Give every edge any positive activity, and put $R=\max(q,q^{-1})$ for $q>0$. The marginal of the random-cluster measure on $G\cup H$, restricted to $E(G)$, has density relative to the measure on $G$ in
$[R^{-(b-1)},R^{b-1}]$.
\end{lemma}
\begin{proof}
For edge configurations $A$ on $G$ and $B$ on $H$, let $\pi,\rho$ be their induced partitions of $S$. Identifying the copies of $S$ gives
\[
k(A\cup B)=k(A)+k(B)-b+\ell(\pi,\rho),\quad
\ell=b-|\pi|-|\rho|+|\pi\vee\rho|.
\]
The quantity $\ell$ is the cycle rank of the bipartite incidence multigraph with one edge for each element of $S$, so $0\le\ell\le b-1$. After summing over $B$, the multiplier of the weight of $A$ therefore lies between $C\min(1,q^{b-1})$ and $C\max(1,q^{b-1})$, where $C>0$ is independent of $A$. Dividing by its mean proves the density bound.
\end{proof}
Sequential attachment of three two-terminal gadgets gives the factor $R^3$. This supplies an applicable corrected principle, not a re-verification of all details of ALR's construction. Our main proofs use exact fan transfer and post integration and do not depend on ALR Lemma~A.1.

\section{Reproduction and assurance record}
Run \texttt{python3 code/verify\_release.py} from the package root. The default performs full reconstruction, including the largest circulant polynomial and the two full-core cubics. The optional \texttt{-{}-quick} mode checks existing exact certificates but does not reconstruct every large object; its receipt records that distinction. Python optimisation is supported: proof-critical checks raise exceptions explicitly rather than relying on \texttt{assert}.

The negative-test suite deliberately corrupts a flow coefficient, graph recipe, witness edge list, new polynomial certificate, and $q=1$ limit sign. It also tests a missing certificate, an unavailable compiler, a malformed historical manuscript table, and an altered current printed coefficient. Every invalid case must exit nonzero, including the specified optimised runs. The immutable supplied v0.2 bundle and referee bundle are retained under \texttt{history/}; the active compatibility code has documented repairs. The historical v0.2 presentation checker covers selected recorded assertions, not every number in the current paper.

The generator \texttt{code/generate\_release\_graph.py} streams a base graph or its full bunkbed. The two smaller base edge lists are included; generating the 36-million-vertex full example is unnecessary for verification. The consolidated execution receipt, claim ledger, response to review, and provenance record distinguish executed checks from written general proofs and unresolved research.

\clearpage
\begin{thebibliography}{99}
\bibitem[ALR(2025)]{ALR} Ayyer, A., Linusson, S., \& Ravichandran, M. (2025).
\emph{The bunkbed problem and the random cluster model}. arXiv:2509.18788, version 1.
\url{https://arxiv.org/abs/2509.18788}
\bibitem[Dong(2015)]{Dong} Dong, F. M. (2015). On zero-free intervals of flow polynomials.
\emph{Journal of Combinatorial Theory, Series B, 111}, 181--200.
Preprint: \url{https://arxiv.org/abs/1403.1916}
\bibitem[FKG(1971)]{FKG} Fortuin, C. M., Kasteleyn, P. W., \& Ginibre, J. (1971).
Correlation inequalities on some partially ordered sets.
\emph{Communications in Mathematical Physics, 22}, 89--103.
\url{https://doi.org/10.1007/BF01651330}
\bibitem[GPZ(2025)]{GPZ} Gladkov, N., Pak, I., \& Zimin, A. (2025).
The bunkbed conjecture is false.
\emph{Proceedings of the National Academy of Sciences, 122}(24), e2420725122.
\url{https://doi.org/10.1073/pnas.2420725122}
\bibitem[GPZ author version(2024)]{GPZauthor} Gladkov, N., Pak, I., \& Zimin, A. (2024, October 8).
\emph{The bunkbed conjecture is false}. Author manuscript, Section~8.5 and footnote~3.
\url{https://www.math.ucla.edu/~gladkovna/Papers/BunkBed6.pdf}
\bibitem[H\"aggstr\"om(2003)]{Haggstrom} H\"aggstr\"om, O. (2003).
Probability on bunkbed graphs. In \emph{Proceedings of FPSAC '03: Formal Power Series and Algebraic Combinatorics}.
\url{https://igm.univ-mlv.fr/~fpsac/FPSAC03/ARTICLES/42.pdf}
\bibitem[Hollom(2025)]{Hollom} Hollom, L. (2025).
The bunkbed conjecture is not robust to generalisation.
\emph{European Journal of Combinatorics, 128}, 104188.
\url{https://doi.org/10.1016/j.ejc.2025.104188}
Preprint: arXiv:2406.01790 (2024).
\url{https://arxiv.org/abs/2406.01790}
\end{thebibliography}
\end{document}
